\documentclass[11pt,a4paper]{article}
\usepackage[T1]{fontenc}
\usepackage[utf8]{inputenc}
\usepackage{lmodern,amsmath,amssymb}
\usepackage[margin=25mm]{geometry}
\usepackage{longtable,booktabs,array,calc}
\usepackage[hidelinks]{hyperref}
\hypersetup{pdftitle={The target box: the strong-coupling gap survives local perturbations, so the renormalizati},pdfauthor={navstokgap research working notes}}
\newcommand{\tightlist}{\setlength{\itemsep}{0pt}\setlength{\parskip}{0pt}}
\setlength{\emergencystretch}{3em}
\setcounter{secnumdepth}{0}
\begin{document}
\hypertarget{the-target-box-the-strong-coupling-gap-survives-local-perturbations-so-the-renormalization-group-has-an-explicit-finish-line}{%
\section{The target box: the strong-coupling gap survives local
perturbations, so the renormalization group has an explicit finish
line}\label{the-target-box-the-strong-coupling-gap-survives-local-perturbations-so-the-renormalization-group-has-an-explicit-finish-line}}

The continuous-time expansion of
\href{kogut-susskind-strong-coupling-explicit.md}{the Kogut--Susskind
note} tolerates a bounded, gauge-invariant, few-link perturbation of the
Hamiltonian and a mild deformation of the electric term. That turns the
strong-coupling result into a \textbf{target}: any renormalization
scheme that steers the effective Hamiltonian into the class
\[\mathcal C(g,\eta,\eta_E,q)=\Big\{H_E'-\tfrac1{g^2}\textstyle\sum_pw_p-\sum_jd_j\ :\
E'(\sigma)\ge(1-\eta_E)\,\varepsilon|S(\sigma)|,\ \ d_j\ \text{gauge invariant on}\le q\ \text{links},\ \ \textstyle\sum_{j\ni\ell}\|d_j\|\le\eta\,\tfrac{\hbar c}{a}\Big\}\]
at some scale \(a\) with \(g^2\) and \(\eta\) in the region below has a
gap of at least \(4\lambda=\frac83g^2(1-\eta_E)(1-\theta)\,\hbar c/a\)
there, and the gap of the original Hamiltonian follows if the scheme is
an exact low-energy reduction. The tolerance for \(SU(3)\), at
\(\theta=1/2\), \(\eta_E=0.1\):

\begin{longtable}[]{@{}lll@{}}
\toprule\noalign{}
\(g^2\) & \(q=4\) (plaquette-like terms) & \(q=8\) (two-plaquette
terms) \\
\midrule\noalign{}
\endhead
\bottomrule\noalign{}
\endlastfoot
\(100\) & \(\eta\le0.29\) & \(\eta\le0.0066\) \\
\(200\) & \(\eta\le0.62\) & \(\eta\le0.036\) \\
\(400\) & \(\eta\le1.3\) & \(\eta\le0.082\) \\
\end{longtable}

in units \(\hbar c/a\), growing linearly in \(g^2\) as
\(\eta_*\simeq g^2/(2e\,2^q)\cdot\frac{1-\eta_E}{3}\). The mass-gap
problem for \(SU(3)\) is therefore the statement that the flow from a
weak bare coupling enters this box. The intermediate region is where
that has to be shown, and the box says exactly what ``enters'' means:
\(g^2\) of order \(10^2\) with few-link corrections of local norm below
a few per cent of \(\hbar c/a\). Constants explicit; the adjacent-growth
count and the quoted tree-graph step are as in the Kogut--Susskind note;
nothing promoted.

\hypertarget{the-perturbed-expansion}{%
\subsection{1. The perturbed expansion}\label{the-perturbed-expansion}}

Take \(H=H_E'-W-D\) with \(W=\frac1{g^2}\sum_pw_p\) as before and
\(D=\sum_jd_j\), each \(d_j\) bounded, gauge invariant, acting on at
most \(q\) links, with the local norm \(\sum_{j\ni\ell}\|d_j\|\le\eta\)
for every link \(\ell\); units \(\hbar c/a=1\). Suppose \(H_E'\) is
diagonal in the character basis with
\(E'(\sigma)\ge\varepsilon'|S(\sigma)|\),
\(\varepsilon'=(1-\eta_E)\varepsilon\).

The Duhamel expansion now has insertions of two kinds. Gauge invariance
of every \(d_j\) keeps all intermediate configurations gauge invariant,
so Gauss's law still forces \(|S_k|\ge4\). The factorization over
disjoint supports holds because each \(d_j\) acts on finitely many
links. A \(d\)-insertion at step \(k\) that touches the current excited
set \(S_{k-1}\) contributes, summed with weights over the terms touching
a given link and over the at most \(2^q\) choices of \(S_k\), at most
\(|S_{k-1}|\,\eta\,2^q\); the energy denominator
\(1/((\varepsilon'-\lambda)|S_{k-1}|)\) cancels the \(|S_{k-1}|\) as
before. The per-step factor of the adjacent-growth count becomes
\[u'=\frac{128Ne/g^2+e\,\eta\,2^q}{\varepsilon'-\lambda},\] and the
first-step factor \(F=8Ne/g^2+e\eta\). The criterion of the
Kogut--Susskind note, Section 2, holds when
\[u'\le\tfrac12\qquad\text{and}\qquad4F\,u'\le\lambda ,\] and then the
gap is at least \(4\lambda\) on the cyclic subspace of local
observables.

\textbf{Proposition.} For \(\lambda=(1-\theta)\varepsilon'\) and
\(u'\le\frac12\), \(4Fu'\le\lambda\), every
\(H\in\mathcal C(g,\eta,\eta_E,q)\) has a unique ground state and a gap
of at least \(\frac83g^2(1-\eta_E)(1-\theta)\,\hbar c/a\), uniformly in
the volume.

\hypertarget{the-numbers}{%
\subsection{2. The numbers}\label{the-numbers}}

For \(SU(3)\), \(\varepsilon=2g^2/3\). With \(\theta=\frac12\) and
\(\eta_E=0.1\), \(\varepsilon'-\lambda=0.3g^2\), so \(u'\le\frac12\)
reads \[e\,\eta\,2^q\ \le\ 0.15\,g^2-\frac{384e}{g^2},\] which is the
table above, and \(4Fu'\le\lambda=0.3g^2\) is far from binding at these
values (\(4Fu'\le2(0.65+2.7\eta)\) at \(g^2=100\)). The tolerance is
linear in \(g^2\) and exponentially small in the locality \(q\) of the
corrections: \textbf{the box is wide for plaquette-like corrections and
narrow for longer-range ones}, which is the quantitative form of
``irrelevant operators must stay small''.

\hypertarget{what-the-box-asks-of-the-intermediate-region}{%
\subsection{3. What the box asks of the intermediate
region}\label{what-the-box-asks-of-the-intermediate-region}}

A renormalization step at scale \(a\) produces an effective Hamiltonian
at \(2a\). Suppose a scheme exists such that

\begin{enumerate}
\def\labelenumi{\arabic{enumi}.}
\tightlist
\item
  each step is an exact low-energy reduction: the spectrum of \(H(a)\)
  below some \(E_c(a)\) coincides with that of \(H_{\rm eff}(2a)\);
\item
  the effective Hamiltonian stays in the class
  \(\mathcal C(g(a),\eta(a),\eta_E(a),q)\) with \(q\) fixed;
\item
  the coupling runs, \(g(2a)>g(a)\), and reaches \(g^2\ge10^2\) after
  finitely many steps with \(\eta,\eta_E\) inside the table.
\end{enumerate}

Then the gap of \(H(a_{\rm UV})\) equals the gap of the effective
Hamiltonian at the final scale, which the Proposition bounds below by
\(\frac83g_s^2(1-\eta_E)(1-\theta)\,\hbar c/a_s\), a positive number in
physical units, uniformly in the volume, and T2\('\) holds along the
trajectory. Items 1--3 in the weak region are the content of the
constructive programme (the Euclidean form of item 2 is Balaban's
small-field effective action); in the intermediate region they are the
open problem. The box makes the finish line explicit: the flow has to
arrive at \(g^2\sim10^2\) with corrections of local norm a few per cent
of \(\hbar c/a\) if they extend over two plaquettes, or of order one if
they are single-plaquette terms.

The one-loop distance from \(g^2=1/2\) to \(g^2=100\) is
\((2-0.01)/0.0966\simeq21\) doublings; the running is perturbative only
at the start of that stretch, and how the corrections behave over the
rest of it is precisely the question.

\hypertarget{consequence-for-state}{%
\subsection{4. Consequence for STATE}\label{consequence-for-state}}

The strong-coupling side is now a forgiving explicit region rather than
a point at infinity, and the mass-gap problem for \(SU(3)\) is the
statement that a renormalization scheme satisfying items 1--3 exists
through the intermediate region. Every quantity in the target is a
number: \(g^2\sim10^2\), corrections of locality \(q\) with local norm
below \(\eta_*(g,q)\) from Section 2, and the resulting gap
\(\frac83g^2(1-\eta_E)(1-\theta)\,\hbar c/a\).
\end{document}
